% Exact frozen candidate proof: one-dimensional endpoint variants only.

\begin{proof}
For \(s\in\{r,R\}\), let \(\eta_s\) denote the \(s\)-Gauss point.
The absolute interior of \(A\) is obtained by deleting precisely the Gauss
point corresponding to each included endpoint.  In particular it contains
the strict radial core
\[
                         A^-_{r,R}=\{r<|T|<R\}
\]
and every nonskeletal residue branch at an included endpoint.

Here are the precise topological inputs.  Over an algebraically closed
complete field, the decomposition of a generalized annulus off its skeleton
into open balls, together with the closure of each ball at its retraction
point, is \cite[Lemma~2.12]{BPR13Structure}; that paper assumes algebraic
closedness throughout.  Over an arbitrary complete field, the corresponding
components are virtual discs by
\cite[\S\S3.6.35.5--3.6.35.10, especially 3.6.35.8]
{Ducros24CurveStructure}, and the skeleton and retraction commute with
complete scalar extension.  Berkovich's relative-interior theorem and
composition formula \cite[Proposition~1.5.5(i)--(ii)]{Berk93}, applied in the
boundaryless analytic torus, identify the relative topological interior with
the absolute analytic interior.  A closed endpoint Gauss point is not in that
interior because the ambient torus skeleton continues through it, whereas
every nonskeletal endpoint branch is ambient open.  This proves the asserted
description of \(\operatorname{Int}(A)\).

Restriction of \eqref{eq:p062-annulus-endpoint-zero-current} to
\(A^-_{r,R}\), followed by the arbitrary-field strict-open annulus theorem,
gives unique \(a,b\in\mathbf R\) such that
\begin{equation}\label{eq:p062-annulus-endpoint-core-formula}
                         h=a\log|T|+b
                         \qquad\text{on }A^-_{r,R}.
\end{equation}
Continuity along the annular skeleton gives the same formula at every
included endpoint Gauss point.  It remains to prove it on the nonskeletal
endpoint branches.

Fix an included endpoint \(s\in\{r,R\}\), and let \(y\in A\) be a rigid
point with \(|T(y)|=s\).  Let
\[
                  P(T)=T^d+\sum_{i=0}^{d-1}c_iT^i
\]
be the monic irreducible polynomial of \(T(y)\) over \(K\).  Completeness of
\(K\) gives a unique extension of its absolute value to every finite
extension.  Hence all roots of \(P\), counted with inseparable multiplicity,
have absolute value \(s\), and the elementary symmetric functions give
\begin{equation}\label{eq:p062-annulus-endpoint-coefficient-bounds}
                  |c_i|\leq s^{d-i},
                  \qquad |c_0|=s^d=:S.
\end{equation}
If \(|T|<s\), the constant term strictly dominates every other term of
\(P\); if \(|T|>s\), the leading term strictly dominates.  Therefore
\begin{equation}\label{eq:p062-annulus-endpoint-open-tube}
 X_y:=\{|P(T)|<S\}\subseteq\{|T|=s\}\subseteq A.
\end{equation}
Moreover \(X_y\) is ambient open in the analytic torus.  Thus it is contained
in the topological, hence absolute analytic, interior of \(A\).  Notice that
no assertion that \(X_y\) lies in one specified residue branch is needed.

The polynomial \(P\) induces a finite flat morphism of degree \(d\)
\begin{equation}\label{eq:p062-annulus-endpoint-open-map}
 p\colon X_y\longrightarrow D^-_S:=\{|Z|<S\},
                  \qquad Z=P(T).
\end{equation}
Indeed, it is the base change of the finite free polynomial morphism, with
basis \(1,T,\ldots,T^{d-1}\).  Let
\[
                         H=\operatorname{Tr}_p(h|_{X_y}),
\]
where the fibre trace uses analytic local degrees.  The finite-flat
continuous-trace theorem and
\eqref{eq:p062-annulus-endpoint-zero-current} make \(H\) continuous with
\(
                    \mathrm d'\mathrm d''[H]=0
\)
on the standard strict open disc \(D^-_S\).  The arbitrary-field strict-open
disc theorem makes \(H\) constant.  The zero fibre is the single analytic
point
\(
                \mathcal M(K[T]/(P))=\{y\}
\)
with local degree \(d\), and hence
\begin{equation}\label{eq:p062-annulus-endpoint-trace-at-zero}
                              H(0)=d\,h(y).
\end{equation}

We next identify the limiting fibre, without defining a closed trace of
\(h\).  The closed polynomial tube is
\begin{equation}\label{eq:p062-annulus-endpoint-closed-tube}
 \overline X_y:=\{|P(T)|\leq S\}=\{|T|\leq s\}.
\end{equation}
The equality follows directly from
\eqref{eq:p062-annulus-endpoint-coefficient-bounds} and the same strict
dominance argument.  The polynomial extends to a finite flat degree-\(d\)
map
\[
             \overline p\colon\overline X_y\longrightarrow
             D^+_S:=\{|Z|\leq S\}.
\]
Only this finite map and its compact source will be used; the function \(h\)
is not extended to \(\overline X_y\setminus A\).

Let \(\xi_S\) be the boundary Gauss point of \(D^+_S\).  Then
\begin{equation}\label{eq:p062-annulus-endpoint-single-gauss-fibre}
       \overline p^{-1}(\xi_S)=\{\eta_s\},
       \qquad \deg_{\eta_s}(\overline p)=d.
\end{equation}
To see this, choose \(u\in K^\times\) with \(|u|=s\), and normalize
\[
                Q(U)=u^{-d}P(uU),
                \qquad V=u^{-d}Z.
\]
All coefficients of \(Q\) are power-bounded, while its leading and constant
coefficients are units.  Its reduction \(q(U)\) therefore has degree exactly
\(d\), and at the source and target Gauss points the residue-field map is
\[
  \widetilde K(V)\longrightarrow\widetilde K(U),
                  \qquad V\longmapsto q(U).
\]
The element \(q(U)\) is transcendental and
\(
  [\widetilde K(U):\widetilde K(q(U))]=d
\), including the inseparable degree.  The graded fundamental inequality
thus gives local degree at least \(d\) at \(\eta_s\).  Finite flatness makes
the sum of the local degrees in the fibre equal to \(d\), proving
\eqref{eq:p062-annulus-endpoint-single-gauss-fibre}.

For \(0<q<S\), let \(\xi_q\) be the \(q\)-Gauss point of \(D^-_S\).  We
claim that, as \(q\) increases to \(S\), the entire finite fibre
\(\overline p^{-1}(\xi_q)\) eventually lies in every prescribed
neighbourhood of \(\eta_s\) in \(\overline X_y\).  If \(W\) is such a
neighbourhood, its compact complement
\(
                  C=\overline X_y\setminus W
\)
has compact, hence closed, image under the finite map \(\overline p\).
By \eqref{eq:p062-annulus-endpoint-single-gauss-fibre}, this image does not
contain \(\xi_S\).  Since \(\xi_q\to\xi_S\), it follows that
\(
                  \overline p^{-1}(\xi_q)\subset W
\)
for all sufficiently large \(q<S\).

Every point above such a \(\xi_q\) satisfies \(|P|=q<S\), and hence belongs
to \(X_y\subset A\), where \(h\) is defined.  Given \(\varepsilon>0\), use
continuity of \(h\) at \(\eta_s\) to choose \(W\) so that
\(
        |h-h(\eta_s)|<\varepsilon
\)
on \(W\cap A\).  Since the total calibrated degree of every fibre is \(d\),
the preceding compact-complement convergence gives
\[
             |H(\xi_q)-d\,h(\eta_s)|\leq d\varepsilon
\]
for all sufficiently large \(q<S\).  Therefore
\begin{equation}\label{eq:p062-annulus-endpoint-trace-limit}
                         H(\xi_q)\longrightarrow d\,h(\eta_s).
\end{equation}
This is a direct endpoint limit.  It neither evaluates \(h\) outside \(A\)
nor invokes a closed trace of \(h\).

The function \(H\) is constant on \(D^-_S\).  Combining
\eqref{eq:p062-annulus-endpoint-trace-at-zero} and
\eqref{eq:p062-annulus-endpoint-trace-limit} yields
\begin{equation}\label{eq:p062-annulus-endpoint-rigid-value}
                              h(y)=h(\eta_s).
\end{equation}
Rigid points are dense in a Hausdorff strict analytic space over a
nontrivially valued field
\cite[Proposition~2.1.15]{Berk12}; applying this locally on \(A\), every
endpoint point is a limit of rigid endpoint points.  Hence continuity,
\eqref{eq:p062-annulus-endpoint-core-formula}, and
\eqref{eq:p062-annulus-endpoint-rigid-value} prove
\eqref{eq:p062-annulus-endpoint-affine} on all of \(A\).

Finally, \(r<R\), so \(\log|T|\) is nonconstant on the annular skeleton.
This proves uniqueness of \(a,b\).
\end{proof}
